\documentclass{article}
\usepackage{hyperref,amsmath}
\title{Preserving a carried effect across restart and supplier loss}
\author{}
\date{26 November 2026}
\begin{document}
\maketitle

\paragraph{Scope.}
This is a positive held-in I1 development envelope, not an evaluation freeze.
It fixes a civil report/station acceptance task, one independent productive
task, one restarted owner, one optional loss of the replacement path after
preparation, and a strict same-obligation deadline. In the replay, each
network request uses the pinned SubT packet implementation. The station
effect, its atomicity and the routes and timings are declared application
conditions. The monitor counts committed station effects and their times,
not merely packet arrivals or acknowledgement. No physical cargo atomicity,
mission generalization, inventory disposition or independent case is claimed.

\paragraph{Matched restart pair.}
The intent contains immutable obligation ID 7. An old supplier either already
has caused the station effect, or its old attempt is terminal without effect.
The old acknowledgement is lost in both. The restarted owner has no reliable
read of the station in either, and no old attempt can act later in this pair.
A fresh owner subprocess receives only the serialized view, plus the installed
policy, while the station and native network persist. The declared branch in
which Q loses its path after preparation then requires use of the backup.
Executing once at a different supplier can still execute the station
effect a second time overall. In the paid history, a blind later acceptance
duplicates. In the terminal unpaid history, withholding misses service.
Any same-view policy that cannot consult or share effect-scoped state and
must settle on a binary execute/withhold decision fails on at least one;
randomizing cannot guarantee both. If a policy chooses to execute with
probability $p$, the two failure probabilities are $p$ and $1-p$.

A truthful application read delivered by cut 2 permits the correct choice
in the terminal pair. This consumes one charged notification byte and
does not imply that the first attempt cannot still act.
In a distinct history the old attempt is still entitled to take effect,
at tick 6; a truthful negative at cut 2 followed by an unconditional backup
causes two effects. Qualified terminal fencing or exclusion must constrain
all competing future effects until the new one commits. A late reading at
tick 5 does not revise the original cut.

\paragraph{Effect-local alternative.}
Keep the consume-once state at the actual station transition, with the same
identity across supplier and attempt generations. Its effect and durable
record are one atomic transition in this model. Each attempted delivery
either establishes the effect and record, or encounters a committed record
and leaves the effect alone. The induction takes logical OR of the initial
effect state and each delivered attempt; only a false-to-true transition
changes the effect count. The live premise is that an authorized eligible
attempt reaches the gate before the deadline. For this declared envelope:
old completion may occur at 1; healthy Q completes at 4; if Q fails at 4
after preparation at 3, the pre-staged backup completes at 6. A still-entitled
old successful attempt at tick 6 is covered by the same atomic boundary.
By contrast, preparing the backup only after seeing Q's loss allows
completion at 8, missing the tick-7 requirement. An independent task commits
at tick 3 in all executions, during acknowledgement loss.
There is an explicit cost: pre-staging the backup spends an extra setup
byte on the no-failure branch, with no extra preparation byte relative to
reactive setup on the failure branch. The gate is not a way to erase
actual packet debits. A separate typed read and the gate can each be
cheaper in different branches, and the read is insufficient for a still-live
old attempt.

\paragraph{Supplier switching without a new signal.}
The reactive backup can derive its choice from a native one-byte Q-to-backup
notice at tick 5, rather than reading the simulator's hidden fault flag.
Erasing that notice isolates the no-information case. When all authorized
attempts use the effect-local gate, a pre-staged timer may instead send
the backup at 6 whether or not Q succeeds. The timer needs no notice,
extends the positive service floor to the joint loss of Q and its cue,
and substitutes an actual backup request for the cue on a healthy branch.
In the tested byte model the two gated policies tie on total transmitted bytes
when Q works; on the failed-but-signaled branch the timer saves the cue byte.
The ledgers and receiver use are not interchangeable, and the trade would
change if a backup attempt cost more than the cue or the two used constrained
different stations. The timer also sends a redundant request on a healthy
branch. It has no exactly-once guarantee with independent supplier caches.

\paragraph{Existing mechanism and replay.}
\href{https://web.stanford.edu/~ouster/cgi-bin/papers/rifl.pdf}{Lee et al.,
RIFL (SOSP 2015), Sections 3 and 4} require the durable completion record
and mutation to be atomic, and require retries to rendezvous with it even
after migration. The common station transition models that strong interface
and ties a same-information factored direct implementation. The supplier-local
case is a boundary ablation rather than a RIFL baseline. The closest-route
or return-to-contact choice by itself does not provide the missing effect
and identity semantics; no comparison with a full navigation stack is made.
A physical task where no reachable controller can implement the atomic
transition needs separate field evidence, an alternative terminal witness,
or a weaker service contract.

The replay compares 540 joint native cases against a direct calculation
including debit and effect counts, and passes their actual terminal events
to the common exactly-once admission interpreter. Ten selected cases use
a freshly launched recipient subprocess while the same native plant runs.
Lean checks the indistinguishable choice, the limited sufficiency of a
negative read, and the gate state/cutoff kernels. The replay intentionally
varies late, pending, missing or absent gate conditions outside the positive
envelope. It does not cover arbitrary other deadlines or second-gate failure.
\end{document}
